\section{Resolving Tail Ambiguity through Revelation}
\label{sec:revelation}
\subsection{Tail identification value}
\label{sec:tiv}
Fix \(f,\Gamma\), and a cutoff \(\eta\in[0,1]\). Throughout this section, expectations are over the target covariate distribution \(P_X\). An audit reveals existing outcomes: its indicator \(A\le1-S\) and intensity \(\lambda(x)=P(A=1\mid X=x,S=0)\) specify the share of originally hidden outcomes recovered at \(x\).

We hold the historical observed law and the audited population contribution fixed; the remaining
hidden outcomes retain the specified sensitivity bounds. Formally, let \(p_A(x)=P(Y=1\mid X=x,A=1)\)
and \(q(x)=P(Y=1\mid X=x,S=0,A=0)\). We impose the original interval \([\underline
p_\Gamma(x),\overline p_\Gamma(x)]\) separately on these audited and still-hidden outcome
probabilities wherever their component has positive mass; \(q\) is otherwise free within it. This
defines the retained-sensitivity information model (Appendix~\ref{app:retained-class}).
Identification treats the audited population contribution as known; finite audited samples require
estimation of that contribution.

The retained lower and upper objectives, \(L_{f,\eta}(\lambda)\) and \(U_{f,\eta}(\lambda)\), share the known contribution and differ only on the remaining hidden mass (Appendix~\ref{app:tiv}). Write \(s_\Gamma=\overline p_\Gamma-\underline p_\Gamma\). Their gap is
\begin{equation}
\begin{aligned}
\tau_{f,\eta}(x)&:=\frac{1-\pi(x)}{\beta}s_\Gamma(x)|d_{f,\eta}(x)|,\\
W_{f,\eta}(\lambda)&:=U_{f,\eta}(\lambda)-L_{f,\eta}(\lambda)\\
&=\mathbb E[(1-\lambda)\tau_{f,\eta}]
=W_{f,\eta}(0)-\mathbb E[\lambda\tau_{f,\eta}].
\end{aligned}
\label{eq:tiv-identity}
\end{equation}
The tail identification value (TIV) density \(\tau_{f,\eta}\) combines three quantities: the population share with hidden outcomes, the allowed range of their outcome probability, and the change in excess loss between the two labels. It measures ambiguity per unit target-population mass. Revealing a fraction \(\lambda\) removes exactly \(\mathbb E[\lambda\tau_{f,\eta}]\), a value calculable before seeing the audited labels. When both labels produce the same excess loss, this value is zero even if the outcome is uncertain.

\subsection{Minimum-cost revelation and certificate precision}
\label{sec:frontier}
Let \(c(x)\ge0\) be an integrable cost density per unit target-population mass, so allocation \(\lambda\) costs \(\mathbb E[c\lambda]\). If one already-hidden label costs \(\kappa(x)\), then \(c=(1-\pi)\kappa\). For a target gap \(\varepsilon\ge0\), the minimum cost is
\begin{equation}
 \operatorname{Cert}_{f,\eta}(\varepsilon)
 :=\inf_{0\le\lambda\le1}
 \{\mathbb E[c\lambda]:W_{f,\eta}(\lambda)\le\varepsilon\}.
\label{eq:cert-frontier}
\end{equation}

\begin{theorem}[Minimum-cost fixed-threshold revelation]
\label{thm:cert-frontier}
Under the retained-sensitivity model, this infimum is attained. An optimum reveals zero-cost positive-TIV locations first and, on \(c>0\), satisfies
\[
 \lambda^\star(x)=
 \begin{cases}
 1,&\tau_{f,\eta}(x)/c(x)>t,\\
 0,&\tau_{f,\eta}(x)/c(x)<t,
 \end{cases}
\]
with fractional revelation on ties, for some \(t\in[0,\infty]\).
\end{theorem}

Thus acquisition follows decreasing ambiguity reduction per unit cost until the target is reached, using a partial terminal tie group when needed. This is the fractional-knapsack rule \citep{dantzig1957discrete}. Appendices~\ref{app:frontier} and~\ref{app:finite-cell} give the proof, boundary cases, and finite-cell computation. Among already-hidden candidates with positive label cost, its operational score is
\[
 \frac{\tau_{f,\eta}(x)}{c(x)}
 =\frac{s_\Gamma(x)|d_{f,\eta}(x)|}{\beta\kappa(x)}.
\]
The hidden-population fraction cancels because the candidate is already known to be hidden.

\label{sec:sequential}
The retained endpoints minimize their cutoffs separately. Their optimized width satisfies
\begin{equation}
 I^{\mathrm{ret}}_{\Gamma,\lambda}(f)
 :=\min_{\eta\in[0,1]}U_{f,\eta}(\lambda)
   -\min_{\eta\in[0,1]}L_{f,\eta}(\lambda)
 \le W_{f,0}(\lambda).
\label{prop:uniform-threshold}
\end{equation}
Theorem~\ref{thm:cert-frontier} at \(\eta=0\) therefore minimizes cost for the sufficient requirement \(W_{f,0}(\lambda)\le\varepsilon\), which guarantees optimized width at most \(\varepsilon\). Appendix~\ref{app:optimized-width} proves this bound and formulates direct optimization of the optimized interval width.

Our sequential audit procedure retrains the predictor after each batch, uses the current
upper-endpoint minimizing cutoff in the retained plug-in objective, and refreshes TIV scores. Its
historical outcome fit remains fixed, preserving the original sensitivity bounds. Audited labels
enter through their observed losses. Candidates are ranked by estimated
\(s_\Gamma|d_{f,\eta}|/(\beta\kappa)\), a batch is revealed, and the procedure repeats.
Figure~\ref{fig:revelation-results} evaluates the fixed and sequential protocols;
Appendix~\ref{app:repayment-audits} specifies fitting and evaluation splits.
